\chapter{Анализ результатов}
\label{ch:analysis}

В таблице~\ref{tab:summary} результаты этапов 4--8 сопоставлены с характеристиками системы из таблицы~\ref{tab:system}.
Так как количество фактических операций в разных режимах различается, для сравнения используются нормализованные метрики: cycles/access, L1d miss rate, число тактов ожидания на промах и branch miss rate.

\begin{table}[H]
\caption{Связь наблюдаемых изменений с характеристиками системы}
\label{tab:summary}
\small
\begin{tabularx}{\textwidth}{|>{\raggedright\arraybackslash\hsize=.8\hsize}X|c|>{\raggedright\arraybackslash\hsize=.7\hsize}X|>{\raggedright\arraybackslash\hsize=1.5\hsize}X|}
	\hline
	Причина & Этап & Характеристика системы & Наблюдаемое изменение \\
	\hline
	Выход рабочего набора за пределы уровней кэша & 4 & L1d 128~KiB, L2 16~MiB & До 64~KiB промахов L1d нет, после 128~KiB появляются промахи L1d (до 2,8\%), после 16~MiB число тактов ожидания на промах растёт с 0,2--0,3 до 0,43--0,47; cycles/access растёт лишь с 1,05 до 1,20 \\
	\hline
	Ухудшение пространственной локальности & 5 & Строка кэша 128~B & cycles/access растёт с 1,15 при stride~=~1 до 18,7 при stride~=~16, когда каждое обращение попадает в новую строку (97\% промахов L1d) \\
	\hline
	Изменение размера элементов & 6 & L1d 128~KiB, L2 16~MiB & При N = 12288 рабочий набор float (96~KiB) помещается в L1d, а double (192~KiB) --- нет: 0,01\% и 1,42\% промахов L1d; при N = 1572864 тактов ожидания на промах 0,27 у float и 0,41 у double \\
	\hline
	Увеличение числа TLB misses & 7 & Страница 16~KiB & Промахи L1 dTLB появляются при переходе от 128 к 256 страницам (cycles/access растёт с 21 до 26), промахи L2 TLB --- при переходе от 2048 к 4096 страницам \\
	\hline
	Увеличение числа ошибок предсказания переходов & 8 & --- & Branch miss rate 40,3\% против практически нулевой, число тактов больше в 8,6 раза, около 17 тактов на одну ошибку \\
	\hline
\end{tabularx}
\end{table}

\section{Выводы}

В результате работы получены аппаратные показатели производительности Apple M3 Pro для последовательного и разреженного доступа к данным и для условных переходов с разной предсказуемостью.
Наибольшее влияние на время выполнения оказывают пространственная локальность и ошибки предсказания переходов.
Выход рабочего набора за пределы L1d и L2 при последовательном доступе увеличивает cycles/access лишь примерно на 15\%, так как аппаратная предвыборка заранее загружает строки кэша.
